\pdfoutput=1
% Options for packages loaded elsewhere
\PassOptionsToPackage{unicode}{hyperref}
\PassOptionsToPackage{hyphens}{url}
\documentclass[
  11pt,
]{article}
\usepackage{xcolor}
\usepackage[margin=1in]{geometry}
\usepackage{amsmath,amssymb}
\setcounter{secnumdepth}{-\maxdimen} % remove section numbering
\usepackage{iftex}
\ifPDFTeX
  \usepackage[T1]{fontenc}
  \usepackage[utf8]{inputenc}
  \usepackage{textcomp} % provide euro and other symbols
\else % if luatex or xetex
  \usepackage{unicode-math} % this also loads fontspec
  \defaultfontfeatures{Scale=MatchLowercase}
  \defaultfontfeatures[\rmfamily]{Ligatures=TeX,Scale=1}
\fi
\usepackage{lmodern}
\ifPDFTeX\else
  % xetex/luatex font selection
\fi
% Use upquote if available, for straight quotes in verbatim environments
\IfFileExists{upquote.sty}{\usepackage{upquote}}{}
\IfFileExists{microtype.sty}{% use microtype if available
  \usepackage[]{microtype}
  \UseMicrotypeSet[protrusion]{basicmath} % disable protrusion for tt fonts
}{}
\makeatletter
\@ifundefined{KOMAClassName}{% if non-KOMA class
  \IfFileExists{parskip.sty}{%
    \usepackage{parskip}
  }{% else
    \setlength{\parindent}{0pt}
    \setlength{\parskip}{6pt plus 2pt minus 1pt}}
}{% if KOMA class
  \KOMAoptions{parskip=half}}
\makeatother
\setlength{\emergencystretch}{3em} % prevent overfull lines
\providecommand{\tightlist}{%
  \setlength{\itemsep}{0pt}\setlength{\parskip}{0pt}}
\usepackage[]{natbib}
\bibliographystyle{plainnat}
\usepackage{orcidlink}
\usepackage{bookmark}
\IfFileExists{xurl.sty}{\usepackage{xurl}}{} % add URL line breaks if available
\urlstyle{same}
\hypersetup{
  pdftitle={Job Security in the Age of AI: A Structured Forecasting Approach},
  pdfauthor={Gergely Máté },
  hidelinks,
  pdfcreator={LaTeX via pandoc}}

\title{Job Security in the Age of AI: A Structured Forecasting Approach}
\author{Gergely Máté \orcidlink{0009-0003-6124-5415}}
\date{February 2026}

\begin{document}
\maketitle
\begin{abstract}
As Artificial Intelligence (AI) capabilities advance in cognitive and
sensory domains, the landscape of global labor markets faces significant
uncertainty. This paper proposes a structured forecasting framework to
estimate job resilience by decomposing occupations into five core
dimensions: cognitive processing, physical execution, social
interaction, sensory perception, and environmental adaptability. By
applying a weighted `Human Advantage Score' and accounting for economic
adoption factors, the model provides a systematic method for quantifying
occupational safety. We demonstrate the utility of this framework
through case studies of construction, software development, and
dentistry, offering a rational basis for career orientation and policy
planning in the AI transition era.
\end{abstract}

\section{Motivation and Scope}\label{motivation-and-scope}

The integration of Artificial Intelligence (AI) into the global economy
represents a paradigm shift in labor dynamics. While previous waves of
automation primarily impacted manual labor, current progress in
generative AI and reasoning models targets cognitive processing---the
act of thinking itself. As cognitive tasks are increasingly delegated to
AI systems, many activities dependent on these skills may face economic
obsolescence as automated alternatives become significantly more
cost-effective than human labor. This transition introduces profound
uncertainty into labor markets, threatening the income stability of
millions of professionals. Systematically analyzing and categorizing
occupations is essential to mitigating transition risks and informing
strategic career orientation.

To address these challenges, this paper proposes a structured
forecasting framework that decomposes occupations into five core
dimensions: Cognitive, Physical, Social, Sensory, and Adaptive. By
applying informed estimates of AI capabilities relative to these skills,
the model calculates an intermediate Job AI Resilience Score based on
comparative human advantages. This metric is further refined by two
critical variables: professional maturity---where mastery and tacit
knowledge provide additional resilience---and socio-economic adoption
factors that influence the actual speed of AI integration.

The ultimate objective of this methodology is to provide a rational,
evidence-based classification of occupations into four Risk Timeline
Categories. By moving beyond speculative fear and into structured
estimation, this framework allows individuals and policymakers to assess
occupational safety with greater precision.

\section{Related Work}\label{related-work}

The systematic assessment of labor displacement by technology has
evolved from broad occupational analysis to more granular, task-based
frameworks. Frey and Osborne (2017) famously identified ``engineering
bottlenecks'' - such as social intelligence and complex perception -
that limit the immediate reach of automation. Similarly, Brynjolfsson et
al.~(2018) proposed a ``Suitability for Machine Learning'' (SML) rubric,
arguing that while few occupations are fully automatable, most consist
of tasks with varying degrees of AI vulnerability.

From an economic perspective, Acemoglu and Restrepo (2019) emphasize
that the risk to labor is driven by the ``displacement effect,'' where
capital-intensive technology replaces human tasks, balanced against a
``productivity effect'' that can create new roles. Current market data
from the World Economic Forum (2025) suggests that as generative AI
matures, high-wage cognitive roles - previously considered safe - are
experiencing accelerated exposure. This paper builds upon these
foundations by proposing a weighted scoring model that combines these
technical bottlenecks with socio-economic adoption factors.

\section{The Methodology}\label{the-methodology}

Any job can be decomposed based on these fundamental dimensions of
required human capabilities:

\begin{itemize}
\tightlist
\item
  \textbf{Cognitive Processing} (thinking, analyzing, planning,
  knowledge application),
\item
  \textbf{Physical Execution} (movement, dexterity, strength),
\item
  \textbf{Social Interaction} (communication, relationships, emotional
  intelligence, persuasion),
\item
  \textbf{Sensory Perception} (seeing, hearing, touching, smelling,
  tasting, balance),
\item
  \textbf{Environment Adaptability} (handling changing conditions,
  environments).
\end{itemize}

Development of AI primarily touches cognitive and sensing capabilities:
AI models excel at applying knowledge and pattern recognition (including
visual and auditory patterns). Thinking models can reason efficiently,
and AIs tend to get better and better at solving academic problems.
Models are also capable of text, audio, image and video generation, and
thus may be capable of solving communication tasks. However, AI driven
robots acting in the physical world are currently at a very early stage.
Although there are some robots capable of moving and manipulating
objects with varying precision, these are still rare, mostly in
prototype phase, and costly. AI also struggles to navigate unseen or
rapidly changing environments, and has edge hardware limitations:
systems showing higher intelligence usually run on energy demanding GPUs
in large clusters of servers, and doesn't fit into autonomous small
robotic bodies. Therefore jobs involving field work are far less
vulnerable to AI compared to office work.

After decomposing the job under scrutiny, we score each dimension based
on how resilient it is to current progress in AI capabilities. This
scoring is inherently heuristic, relying on informed estimations related
to the specific activities the job requires.

Then we can apply the weighted job vulnerability formula, where higher
scores indicate greater resilience against AI disruption:

\[Job\ AI\ Resilience\ Score = \sum_{i} (Component\ Weight_i \times Human\ Advantage\ Score_i).\]

\clearpage

Humans advance their art during their years in a profession. To account
for tacit knowledge, we distinguish among beginner, advanced and master
level based on experience by applying -1, 0 and +1 modifiers. With this
we try to capture the fact that jobs of more skilled, experienced
professionals are less vulnerable to AI than jobs at the beginner level.

We use this score to classify jobs into risk timeline categories of
\textbf{immediate risk}, \textbf{medium term risk}, \textbf{long term
safety} and \textbf{very low risk} (probably indefinite safety).

Finally, we modify the result based on four important adoption factors:
first, for companies there is a higher incentive to replace workers with
higher wages to cut costs, second, for some activities replacing humans
with machines can have non-economical benefits - for example reducing
risk of harm -, third, some development (particularly robots) can be
prohibitively costly, and fourth, there are jobs (especially where
something is ``performed'' in a broad sense) where most humans
inherently prefer to see humans.

\section{Application of the Analysis
Framework}\label{application-of-the-analysis-framework}

\subsection{1. Job Composition (Total
100\%)}\label{job-composition-total-100}

Quickly estimate the percentage breakdown of:

\begin{itemize}
\tightlist
\item
  \textbf{Thinking} (analysis, planning, creativity): \_\_\_\_\_\%
\item
  \textbf{Physical} (movement, dexterity, strength): \_\_\_\_\_\%
\item
  \textbf{Social} (communication, emotional intelligence): \_\_\_\_\_\%
\item
  \textbf{Perception} (seeing, hearing, feeling): \_\_\_\_\_\%
\item
  \textbf{Adaptation} (handling changing environments): \_\_\_\_\_\%
\end{itemize}

\subsection{2. Human Advantage Score (1-10 for each
component)}\label{human-advantage-score-1-10-for-each-component}

For each component, take into account the job specific activities, and
score how much advantage humans have:

\begin{itemize}
\tightlist
\item
  \textbf{1-3}: AI already good or improving rapidly
\item
  \textbf{4-6}: AI partially capable but significant gaps
\item
  \textbf{7-10}: AI struggles significantly
\end{itemize}

\subsection{3. Calculate Job Resilience
Score}\label{calculate-job-resilience-score}

\begin{itemize}
\tightlist
\item
  Multiply each component's percentage by its Human Advantage Score
\item
  Add them together for a weighted total out of 10
\item
  Apply experience modifier: Beginner (-1), Advanced (0), Master (+1)
\end{itemize}

\subsection{4. Interpret Results}\label{interpret-results}

\begin{itemize}
\tightlist
\item
  \textbf{Below 4}: High Risk (machines are overtaking the job in 1-5
  years)
\item
  \textbf{4-6}: Medium Risk (5-10 years)
\item
  \textbf{Above 6}: Low Risk (10+ years)
\item
  \textbf{Above 8}: Very Low Risk (15+ years or indefinite human
  advantage)
\end{itemize}

\subsection{5. Adoption Factors Check}\label{adoption-factors-check}

Ask these questions:

\begin{enumerate}
\def\labelenumi{\arabic{enumi}.}
\tightlist
\item
  Is the job's compound annual salary significantly higher than
  potential AI or robot implementation costs?
\item
  Would replacing humans provide substantial benefits beyond cost
  savings?
\item
  Would the implementation cost of an AI or robot solution be
  prohibitively expensive for the foreseeable future?
\item
  Does the job involve performance or service where human authenticity
  itself is inherently valued?
\end{enumerate}

If any of the answer to questions 1 or 2 is ``yes,'' move the job up one
risk level.\\
If any of the answer to questions 3 or 4 is ``yes,'' move the job down
one risk level.

\section{Three Examples}\label{three-examples}

\subsection{Construction Worker}\label{construction-worker}

\textbf{1. Job Composition}

\begin{itemize}
\tightlist
\item
  Thinking: 15\% (problem-solving, basic planning)
\item
  Physical: 50\% (lifting, operating machinery, manual labor)
\item
  Social: 10\% (team coordination)
\item
  Perception: 15\% (spatial awareness, hazard detection)
\item
  Adaptation: 10\% (responding to changing site conditions)
\end{itemize}

\textbf{2. Human Advantage Score}

\begin{itemize}
\tightlist
\item
  Thinking: 4/10 (AI can handle basic plans but struggles with on-site
  decisions)
\item
  Physical: 8/10 (robots still struggle with construction-site mobility
  and strength)
\item
  Social: 6/10 (limited real-time team coordination capabilities)
\item
  Perception: 7/10 (computer vision struggles in unstructured
  environments)
\item
  Adaptation: 9/10 (significant challenges with unpredictable
  environments)
\end{itemize}

\textbf{3. Job Resilience Score}

\begin{itemize}
\tightlist
\item
  Total: 7.15/10
\item
  Beginner: 6.15 (7.15-1)
\item
  Advanced: 7.15
\item
  Master: 8.15 (7.15+1)
\end{itemize}

\textbf{4. Result: Low to Very Low Risk} (10-15+ years)

\textbf{5. Adoption Factors Check:} Safety benefits beyond cost savings
might increase risk level slightly, but construction work still remains
in the low risk category.

\subsection{Software Developer}\label{software-developer}

\textbf{1. Job Composition}

\begin{itemize}
\tightlist
\item
  Thinking: 60\% (coding, problem-solving, design)
\item
  Physical: 5\% (typing, basic computer interaction)
\item
  Social: 20\% (stakeholder communication, requirements gathering)
\item
  Perception: 10\% (code review, testing)
\item
  Adaptation: 5\% (handling changing requirements)
\end{itemize}

\textbf{2. Human Advantage Score}

\begin{itemize}
\tightlist
\item
  Thinking: 4/10 (good at code generation but struggles with
  larger-scale planning)
\item
  Physical: 1/10 (AI doesn't need physical interaction to code)
\item
  Social: 5/10 (struggles with nuanced stakeholder management)
\item
  Perception: 4/10 (effective at code analysis)
\item
  Adaptation: 4/10 (can adapt to changes but misses context)
\end{itemize}

\textbf{3. Job Resilience Score}

\begin{itemize}
\tightlist
\item
  Total: 4.05/10
\item
  Beginner: 3.05 (4.05-1)
\item
  Advanced: 4.05
\item
  Master: 5.05 (4.05+1)
\end{itemize}

\textbf{4. Result:}

\begin{itemize}
\tightlist
\item
  Beginner: High Risk (1-5 years)
\item
  Advanced: Medium Risk (5-10 years)
\item
  Master: Medium Risk (5-10 years)
\end{itemize}

\textbf{5. Adoption Factors Check: }High salaries justify AI
replacement, potentially moving advanced developers toward High Risk,
while master developers with stronger stakeholder relationships remain
in Medium Risk\textbf{.}

\subsection{Dentist}\label{dentist}

\textbf{1. Job Composition}

\begin{itemize}
\tightlist
\item
  Thinking: 25\% (diagnosis, treatment planning, medical knowledge)
\item
  Physical: 30\% (precise hand movements, tool manipulation)
\item
  Social: 15\% (patient communication, anxiety management)
\item
  Perception: 20\% (visual examination, tactile feedback)
\item
  Adaptation: 10\% (responding to patient movement, unexpected findings)
\end{itemize}

\textbf{2. Human Advantage Score}

\begin{itemize}
\tightlist
\item
  Thinking: 5/10 (AI can assist with diagnosis but struggles with
  comprehensive treatment planning)
\item
  Physical: 9/10 (extremely difficult to replicate precise hand
  movements in small, sensitive spaces)
\item
  Social: 7/10 (difficult to replicate bedside manner and real-time
  emotional reassurance)
\item
  Perception: 8/10 (complex integration of visual and tactile feedback
  in variable environments)
\item
  Adaptation: 9/10 (significant challenges adapting to unexpected
  patient movements and oral conditions)
\end{itemize}

\textbf{3. Job Resilience Score}

\begin{itemize}
\tightlist
\item
  Total: 7.5/10
\item
  Beginner: 6.5 (7.5-1)
\item
  Advanced: 7.5
\item
  Master: 8.5 (7.5+1)
\end{itemize}

\textbf{4. Result:}

\begin{itemize}
\tightlist
\item
  Beginner: Low Risk (10+ years)
\item
  Advanced: Low Risk (10+ years)
\item
  Master: Very Low Risk (15+ years or indefinite human advantage)
\end{itemize}

\textbf{5. Adoption Factors Check:} High salaries might justify
automation investment (risk up), but implementation costs for robotic
dentistry would be extremely high, and many patients are likely to
prefer humans (risk down). These factors cancel each other out, keeping
dentistry in the Low to Very Low Risk category.

\section{Conclusion and Future
Outlook}\label{conclusion-and-future-outlook}

This framework provides an immediate, structured, systematic method to
transition from speculative anxiety to rational assessment of
occupational risk. By decomposing occupations into core dimensions,
applying weighted resilience scores, and accounting for economic and
social adoption factors, individuals and policymakers can evaluate job
security in real-time.

However, as a forecasting model, this tool is inherently iterative.
While current resilience scores are based on informed estimates of
existing AI and robotic capabilities, the rapid pace of technological
development necessitates continuous updates to the ``Human Advantage''
metrics. There remains significant room for methodological refinement,
empirical validation, and the integration of shifting socio-economic
adoption factors.

This methodology provides a systematic starting point for individuals
and policymakers to assess occupational safety. I encourage readers to
apply this framework to their own professional trajectories and
contribute to the ongoing refinement of these estimations as we navigate
this transition.

\section{Acknowledgments}\label{acknowledgments}

The author acknowledges the use of various Claude and Gemini models for
structural brainstorming and linguistic refinement.

\section{References}\label{references}

\begin{enumerate}
\def\labelenumi{\arabic{enumi}.}
\tightlist
\item
  Acemoglu, D., \& Restrepo, P. (2019). ``Automation and New Tasks: How
  Technology Displaces and Reinstates Labor.'' Journal of Economic
  Perspectives, 33(2), 3--30.
\item
  Brynjolfsson, E., \& Mitchell, T. (2017). ``What can machine learning
  do? Workforce implications'' \emph{Science}, 358(6370), 1530-1534.
\item
  Frey, C. B., \& Osborne, M. A. (2017). ``The future of employment: How
  susceptible are jobs to computerisation?'' \emph{Technological
  Forecasting and Social Change}, 114, 254-280.
\item
  World Economic Forum. (2025). ``The Future of Jobs Report 2025.''
\end{enumerate}

\end{document}
